% !PRINTED 521
\weberpage{521}
%<<CONT
and which has no element in common with either $A$ or $B$. One
continues in this way until all elements are exhausted, and comes to
the following theorem:

\gesperrt{2. A transitive group has, only when it is imprimitive,
an intransitive normal divisor different
from the identical group. The systems of
intransitivity of the normal divisor are
systems of imprimitivity of the given
group.} % The German source prints "Imprimivit\"at" for "Imprimitivit\"at"; the intended technical term is translated here.

With respect to the equations, having regard to 1.,
we can also express theorem 2. as follows:

\gesperrt{3. An irreducible equation $f(x) = 0$ splits,
when it becomes reducible through the adjunction of all
the roots of a resolvent or of one root of a normal equation
of degree $j$, into several
irreducible factors of equal degree and
of equal group. The number of these factors
is a divisor of the degree $j$ and is equal to $j$
when $j$ is a prime.}

If, then, an irreducible primitive equation is reduced by the
adjunction of the roots of a resolvent, it splits
into linear factors, i.\,e. it is completely solved.
This case always occurs when the degree of the equation
$f(x) = 0$ is a prime, because such an equation cannot
be imprimitive.

\begin{center}
\rule{3cm}{0.4pt}
\end{center}
